\documentclass[aspectratio=169]{beamer}

% ─────────────────────────────────────────────
%  Theme & colours
% ─────────────────────────────────────────────
\usetheme{Madrid}
\usecolortheme{default}
\definecolor{WUBlue}{RGB}{0,62,120}
\definecolor{WUGold}{RGB}{192,145,0}
\definecolor{LightBlue}{RGB}{55,138,221}
\definecolor{DarkBlue}{RGB}{21,47,95}
\definecolor{AccentOrange}{RGB}{224,92,42}

\setbeamercolor{palette primary}{bg=WUBlue,fg=white}
\setbeamercolor{palette secondary}{bg=WUBlue!80,fg=white}
\setbeamercolor{palette tertiary}{bg=WUBlue!60,fg=white}
\setbeamercolor{palette quaternary}{bg=WUBlue,fg=white}
\setbeamercolor{structure}{fg=WUBlue}
\setbeamercolor{section in toc}{fg=WUBlue}
\setbeamercolor{alerted text}{fg=AccentOrange}
\setbeamercolor{block title}{bg=WUBlue,fg=white}
\setbeamercolor{block body}{bg=WUBlue!8}

\setbeamertemplate{navigation symbols}{}
\setbeamertemplate{footline}[frame number]

% ─────────────────────────────────────────────
%  Packages
% ─────────────────────────────────────────────
\usepackage{booktabs}
\usepackage{array}
\usepackage{graphicx}
\usepackage{amsmath}
\usepackage{xcolor}
\usepackage{colortbl}
\usepackage{tikz}
\usepackage{multirow}
\usepackage{tabularx}
\usepackage{adjustbox}
\usepackage{caption}

% ─────────────────────────────────────────────
%  Title page
% ─────────────────────────────────────────────
\title[GEP Index — Interim Results]%
      {\textbf{Geoeconomic Pressure Index}\\[4pt]
       \large Interim Results: Construction, Revision \& Extension to Europe}
\author[Caterina Piacentini]{Caterina Piacentini}
\institute[WU Vienna]{Vienna University of Economics and Business\\
           \smallskip\small Thesis Advisor: Dangl \& Salbrechter methodology}
\date{April 2026}

% ─────────────────────────────────────────────
\begin{document}
% ─────────────────────────────────────────────

\begin{frame}
  \titlepage
\end{frame}

% ─── Outline ─────────────────────────────────
\begin{frame}{Outline}
  \tableofcontents
\end{frame}

% ════════════════════════════════════════════════════════
\section{Index 8 — Baseline Construction}
% ════════════════════════════════════════════════════════

% ─── Slide: what is index_8 ──────────────────
\begin{frame}{GTM-8 Baseline: Eight Geoeconomic Sub-topics}

The \textbf{Guided Topic Model (GTM-8)} expands eight seed clusters
in a 64-dimensional Word2Vec embedding space trained on
30 years of Reuters newswire (1996--2025).

\medskip
\begin{columns}[T]
\column{0.48\textwidth}
\begin{block}{Sub-topics (8)}
\begin{enumerate}\small
  \item Sanctions
  \item Trade War
  \item Export Controls
  \item Tariffs
  \item Financial Coercion
  \item Retaliation
  \item Protectionism
  \item Embargo
\end{enumerate}
\end{block}

\column{0.48\textwidth}
\begin{block}{Key parameters}
\begin{itemize}\small
  \item Cluster size: 100 words each
  \item Gravity: 1.5
  \item $\alpha_\text{max}$: 2.0 rad
  \item FAISS $k$-NN: 5\,000 candidates
  \item Final dictionary: \textbf{493 words}
\end{itemize}
\end{block}
\end{columns}

\end{frame}

% ─── Slide: seeds table index_8 ──────────────
\begin{frame}{Index 8 — Seed Words per Sub-topic}

\begin{center}
\footnotesize
\renewcommand{\arraystretch}{1.2}
\resizebox{\textwidth}{!}{%
\begin{tabular}{@{}llll@{}}
\toprule
\textbf{Sub-topic} & \textbf{Positive seeds} & \textbf{Negative seeds} \\
\midrule
\rowcolor{WUBlue!10}
Sanctions
  & \texttt{economic\_sanctions}, \texttt{sanctions\_regime}
  & \texttt{sanctions\_relief}, \texttt{sanctions\_waiver} \\
Trade War
  & \texttt{trade\_war}, \texttt{trade\_conflict}
  & \texttt{trade\_deal}, \texttt{trade\_pact} \\
\rowcolor{WUBlue!10}
Export Controls
  & \texttt{export\_ban}, \texttt{export\_restriction}
  & \texttt{export\_license} \\
Tariffs
  & \texttt{import\_tariff}, \texttt{customs\_duty}
  & \texttt{free\_trade}, \texttt{tariff\_exemption} \\
\rowcolor{WUBlue!10}
Financial Coercion
  & \texttt{asset\_freeze}, \texttt{frozen\_assets}
  & \texttt{debt\_relief} \\
Retaliation
  & \texttt{retaliation}, \texttt{retaliatory\_measures}
  & \texttt{concessions}, \texttt{goodwill\_gesture} \\
\rowcolor{WUBlue!10}
Protectionism
  & \texttt{protectionism}, \texttt{import\_quota}
  & \texttt{trade\_liberalization} \\
Embargo
  & \texttt{trade\_embargo}, \texttt{oil\_embargo}
  & \texttt{lift\_embargo} \\
\bottomrule
\end{tabular}}
\end{center}

\vspace{2pt}
\footnotesize
Positive seeds \emph{anchor} the subspace; negative seeds act as
\emph{repellors}, pushing the expansion away from resolution/relief language.

\end{frame}

% ═══════════════════════════════════════════════════════
\section{Why We Revised Index 8}
% ═══════════════════════════════════════════════════════

\begin{frame}{Motivation for Revision}

Inspection of the baseline index revealed three issues:

\medskip
\begin{enumerate}
  \item \textbf{Weak positive seeds in key topics.}\\
        \emph{Sanctions:} \texttt{economic\_sanctions} is broad; it matched
        \emph{relief} contexts (debt relief, sanctions waiver) because the
        positive subspace captured generic ``sanctions'' rather than
        \emph{coercive} sanctions.\\[4pt]
        \emph{Financial Coercion:} \texttt{asset\_freeze} pointed to generic
        freezing (court orders, estate law) — not sovereign-level financial
        weaponisation.

  \medskip
  \item \textbf{Export Controls under-specified.}\\
        \texttt{export\_ban} pulled semiconductor-agnostic words;
        the modern strategic-technology dimension (entity lists, chip controls)
        was not captured.

  \medskip
  \item \textbf{Protectionism seeds too narrow.}\\
        \texttt{import\_quota} biased the topic toward quota-specific WTO
        language; the broader economic nationalism framing was missed.
\end{enumerate}

\end{frame}

% ─── Revised seeds ───────────────────────────
\begin{frame}{Index 8 Revised — Updated Seed Words}

\begin{center}
\footnotesize
\renewcommand{\arraystretch}{1.2}
\resizebox{\textwidth}{!}{%
\begin{tabular}{@{}p{3.0cm}p{4.8cm}p{5.0cm}@{}}
\toprule
\textbf{Sub-topic} & \textbf{Positive seeds (revised)} & \textbf{Change vs.\ baseline} \\
\midrule
\rowcolor{WUBlue!10}
Sanctions
  & \texttt{sectoral\_sanctions}, \texttt{targeted\_sanctions}
  & More specific — coercive intent \\
Trade War
  & \texttt{trade\_war}, \texttt{trade\_frictions}
  & \texttt{trade\_conflict}$\to$\texttt{trade\_frictions};
    \texttt{financial\_markets} to negatives \\
\rowcolor{WUBlue!10}
Export Controls
  & \texttt{entity\_list}, \texttt{export\_controls}
  & Captures US BIS / chip controls \\
Tariffs
  & \textit{unchanged}
  & --- \\
\rowcolor{WUBlue!10}
Financial Coercion
  & \texttt{frozen\_russian}, \texttt{payment\_freeze}
  & Sovereign-level; \texttt{debt\_restructuring} to negatives \\
Retaliation
  & \textit{unchanged}
  & --- \\
\rowcolor{WUBlue!10}
Protectionism
  & \texttt{protectionism}, \texttt{economic\_nationalism}
  & \texttt{import\_quota}$\to$\texttt{economic\_nationalism} \\
Embargo
  & \textit{unchanged}
  & --- \\
\bottomrule
\end{tabular}}
\end{center}

\end{frame}

% ═══════════════════════════════════════════════════════
\section{Word Clouds — Aggregated Topics}
% ═══════════════════════════════════════════════════════

\begin{frame}{Dictionary Aggregation Mechanism}

Each sub-topic produces a ranked word list with \emph{projection norms}
as importance weights. The global dictionary is built in two steps
(Dangl \& Salbrechter, Eq.~10):

\medskip
\textbf{Step 1 — Local normalisation (per topic $q$):}
\[
  \tilde{w}_{q,j} = \frac{w_{q,j}}{\max_j\, w_{q,j}}
  \;\in\; [0,1]
\]

\textbf{Step 2 — Global average (across $Q=8$ topics):}
\[
  \bar{w}_j = \frac{1}{Q}\sum_{q=1}^{Q} \tilde{w}_{q,j}
\]

\medskip
Words appearing in \emph{multiple} sub-topics receive proportionally
higher global weights — they are semantically central across
the geoeconomic pressure space.

\medskip
\begin{block}{Result}
  \textbf{493-word dictionary} (US baseline). Top words:
  \texttt{protectionist} (0.409), \texttt{import\_tariffs} (0.365),
  \texttt{impose\_tariffs} (0.340), \texttt{retaliatory\_tariffs} (0.335),
  \texttt{trade\_war} (0.299).
\end{block}

\end{frame}

% ─── Combined word-cloud grid ─────────────────
\begin{frame}{Word Clouds — All 8 Sub-topics (US Baseline)}
  \begin{center}
    \includegraphics[width=0.85\textwidth]%
      {GTM_different_versions/GTM8_results/Combined_Geoeconomic_Pressure_Grid_8.png}
  \end{center}
\end{frame}

% ─── Revised combined grid ────────────────────
\begin{frame}{Word Clouds — All 8 Sub-topics (Revised)}
  \begin{center}
    \includegraphics[width=0.85\textwidth]%
      {GTM_different_versions/GTM8_revised_results/Combined_Geoeconomic_Pressure_Grid_8.png}
  \end{center}
\end{frame}

% ═══════════════════════════════════════════════════════
\section{GEP Index Results — Revised}
% ═══════════════════════════════════════════════════════

\begin{frame}{GEP Monthly Index — Revised (1996–2025)}
  \begin{center}
    \includegraphics[width=0.95\textwidth]%
      {INDEX/index_8_revised/GEP_Monthly_Index_Revised.png}
  \end{center}
  \vspace{-4pt}
  \footnotesize Score $\times 10^{-4}$; article-weighted monthly average.
  Annotated peaks correspond to key geoeconomic episodes.
\end{frame}

\begin{frame}{GEP Daily Index — Revised (1996–2025)}
  \begin{center}
    \includegraphics[height=0.78\textheight]%
      {INDEX/index_8_revised/GEP_Daily_Events_Revised.png}
  \end{center}
\end{frame}

% ═══════════════════════════════════════════════════════
\section{Robustness Checks}
% ═══════════════════════════════════════════════════════

\begin{frame}{Robustness Check — Article- vs.\ Day-Weighted Aggregation}

\textbf{Question:} Does the choice of monthly aggregation method drive the index?

\medskip
Two variants compared:
\begin{itemize}
  \item \textcolor{LightBlue}{\textbf{Baseline}} — article-weighted average:
        $\displaystyle \text{GEP}_m = \frac{\sum_{t\in m} s_t \cdot n_t}{\sum_{t\in m} n_t}$
  \item \textcolor{AccentOrange}{\textbf{Robustness}} — day-weighted average:
        each trading day receives equal weight regardless of article volume
\end{itemize}

\medskip
  \begin{center}
    \includegraphics[width=0.92\textwidth]%
      {INDEX/index_8_revised/robustness_check_GEP_revised.png}
  \end{center}

\vspace{-4pt}
\footnotesize
Correlation $r > 0.95$ confirms that the aggregation choice is \emph{not}
the driving force behind the index shape.

\end{frame}

% ─── GEP vs GPR ──────────────────────────────
\begin{frame}{GEP vs.\ GPR (Caldara \& Iacoviello 2022) — Revised}

\textbf{External benchmark:} the Geopolitical Risk Index (GPRD) captures
\emph{geopolitical} risk; our GEP captures \emph{geoeconomic} pressure.
Both series are z-scored before comparison.

  \begin{center}
    \includegraphics[width=0.92\textwidth]%
      {INDEX/index_8_revised/gep_vs_gprd_timeseries_revised.png}
  \end{center}

\vspace{-4pt}
\footnotesize
Overall Pearson $r \approx 0.13$.
GEP Granger-causes GPRD from lag~2 onwards ($p < 0.05$) —
suggesting geoeconomic pressure \emph{precedes} geopolitical risk spikes.

\end{frame}

\begin{frame}{Cross-Correlation: GEP Revised vs.\ GPRD}

  \begin{center}
    \includegraphics[width=0.88\textwidth]%
      {INDEX/index_8_revised/gep_vs_gprd_crosscorr_revised.png}
  \end{center}

\medskip\small
\begin{itemize}
  \item \textcolor{LightBlue}{\textbf{Blue bars}} (lag $\ge 0$): GEP leads GPR
  \item \textcolor{AccentOrange}{\textbf{Orange bars}} (lag $< 0$): GPR leads GEP
  \item Peak correlation at \emph{positive} lag confirms that geoeconomic
        pressure news tends to \emph{precede} escalation in geopolitical risk
        readings.
\end{itemize}

\end{frame}

% ═══════════════════════════════════════════════════════
\section{Europe GEP Index}
% ═══════════════════════════════════════════════════════

\begin{frame}{Europe GEP Index — Setup}

A parallel GEP index is constructed from a \textbf{Europe-filtered Reuters corpus}
(articles with EU-country subject codes), using an \emph{independently trained}
GTM dictionary.

\medskip
\begin{columns}[T]
\column{0.48\textwidth}
\begin{block}{Corpus statistics}
\small
\begin{tabular}{lrr}
\toprule
 & \textbf{US} & \textbf{EU} \\
\midrule
Avg art./month & 62\,000 & 177\,000 \\
Months & 357 & 348 \\
Dict.\ words & 493 & 657 \\
\bottomrule
\end{tabular}
\end{block}

\column{0.48\textwidth}
\begin{block}{Distribution of GEP$_m$}
\small
\begin{tabular}{lcc}
\toprule
 & \textbf{US} & \textbf{EU} \\
\midrule
Mean & 0.000563 & 0.000400 \\
Std  & 0.000146 & 0.000060 \\
CV   & 0.260    & 0.149    \\
\bottomrule
\end{tabular}
\smallskip\\
EU index is notably \emph{flatter} (lower CV).
This reflects larger article volume causing statistical smoothing.
\end{block}
\end{columns}

\medskip\footnotesize
\textbf{Normalisation:} EPU-style — $\text{GEP\_norm}_m = \dfrac{\text{GEP}_m}{\overline{\text{GEP}}} \times 100$
(long-run mean = 100 for each series independently).

\end{frame}

% ─── Europe seeds vs US seeds ────────────────
\begin{frame}{Europe GTM Seeds — Comparison with US (1/2)}

\begin{center}
\scriptsize
\renewcommand{\arraystretch}{1.15}
\resizebox{\textwidth}{!}{%
\begin{tabular}{@{}p{2.8cm}p{4.2cm}p{5.5cm}@{}}
\toprule
\textbf{Sub-topic}
  & \textbf{US (revised) seeds}
  & \textbf{Europe seeds / key difference} \\
\midrule
\rowcolor{WUBlue!10}
Sanctions
  & \texttt{sectoral\_sanctions}, \texttt{targeted\_sanctions}
  & \texttt{targeted\_sanctions}, \texttt{sanctions\_regime}
    \newline\emph{Retains regime framing — EU sanctions language} \\
Trade War
  & \texttt{trade\_war}, \texttt{trade\_frictions}
  & \texttt{trade\_war}, \texttt{trade\_friction} (singular)
    \newline\emph{Fewer negative seeds (no financial\_markets)} \\
\rowcolor{WUBlue!10}
Export Controls
  & \texttt{entity\_list}, \texttt{export\_controls}
  & \texttt{export\_ban}, \texttt{export\_controls}
    \newline\emph{No entity\_list — absent from EU corpus} \\
Tariffs
  & \texttt{import\_tariff}, \texttt{customs\_duty}
  & same seeds; no \texttt{free\_trade} negative \\
\bottomrule
\end{tabular}}
\end{center}
\end{frame}

\begin{frame}{Europe GTM Seeds — Comparison with US (2/2)}

\begin{center}
\scriptsize
\renewcommand{\arraystretch}{1.15}
\resizebox{\textwidth}{!}{%
\begin{tabular}{@{}p{2.8cm}p{4.2cm}p{5.5cm}@{}}
\toprule
\textbf{Sub-topic}
  & \textbf{US (revised) seeds}
  & \textbf{Europe seeds / key difference} \\
\midrule
\rowcolor{WUBlue!10}
Financial Coercion
  & \texttt{frozen\_russian}, \texttt{payment\_freeze}
  & \texttt{asset\_freeze}, \texttt{frozen\_assets}
    \newline\emph{Generic freeze; \texttt{debt\_forgiveness} as negative} \\
Retaliation
  & \texttt{retaliation}, \texttt{retaliatory\_measures}
  & same seeds \\
\rowcolor{WUBlue!10}
Protectionism
  & \texttt{protectionism}, \texttt{economic\_nationalism}
  & \texttt{protectionism}, \texttt{import\_quota}
    \newline\emph{Quota framing common in EU WTO discourse} \\
Embargo
  & \texttt{trade\_embargo}, \texttt{oil\_embargo}
  & \texttt{trade\_embargo}, \texttt{arms\_embargo}
    \newline\emph{Arms embargo more salient in EU corpus (Russia, Iran)} \\
\bottomrule
\end{tabular}}
\end{center}
\end{frame}

% ─── Europe word clouds ───────────────────────
\begin{frame}{Word Clouds — All 8 Sub-topics (Europe)}
  \begin{center}
    \includegraphics[width=0.85\textwidth]%
      {GTM_different_versions/GTM8_europe_results/Combined_Geoeconomic_Pressure_Grid_8.png}
  \end{center}
\end{frame}

% ─── Europe monthly index ────────────────────
\begin{frame}{GEP Monthly Index — Europe (1996–2025)}
  \begin{center}
    \includegraphics[width=0.95\textwidth]%
      {INDEX/index_8_europe/GEP_Monthly_Index_Europe.png}
  \end{center}
  \vspace{-4pt}
  \footnotesize
  EPU-normalised; long-run mean = 100.
  Lower coefficient of variation (0.149 vs 0.260 US) reflects larger article
  pool and structural dictionary differences.
\end{frame}

% ─── US vs EU comparison ─────────────────────
\begin{frame}{US vs.\ EU GEP — Head-to-Head Comparison}
  \begin{center}
    \includegraphics[width=0.95\textwidth]%
      {INDEX/us_eu_comparison/GEP_Monthly_US_vs_EU.png}
  \end{center}

  \vspace{-2pt}
  \footnotesize
  Both series normalised to mean = 100.
  \textbf{Only 27\% dictionary overlap} — the two indices measure
  \emph{structurally different} facets of geoeconomic pressure and
  cannot be compared at the level of raw scores.
  EU index is notably flatter; absolute level comparisons should be avoided.
\end{frame}

% ─── Dictionary divergence ────────────────────
\begin{frame}{US vs.\ EU Dictionary Divergence}

\begin{columns}[T]
\column{0.45\textwidth}
\begin{block}{Overlap statistics}
\small
\begin{tabular}{lcc}
\toprule
 & \textbf{US} & \textbf{EU} \\
\midrule
Total words & 493 & 657 \\
Common words & \multicolumn{2}{c}{195 (27\%)} \\
Unique & 298 & 462 \\
Mean weight & 0.064 & 0.072 \\
\bottomrule
\end{tabular}
\end{block}

\column{0.50\textwidth}
\begin{block}{Key weight differences}
\small
\begin{tabular}{lc}
\toprule
\textbf{Word} & \textbf{EU $-$ US $\Delta$weight} \\
\midrule
\texttt{import\_duty}       & $+0.230$ \\
\texttt{import\_restrictions} & $+0.200$ \\
\texttt{imposing\_tariffs}  & $+0.160$ \\
\texttt{protectionist}      & $-0.090$ \\
\texttt{trade\_barriers}    & $-0.130$ \\
\texttt{retaliatory\_tariffs} & $-0.090$ \\
\bottomrule
\end{tabular}
\end{block}
\end{columns}

\bigskip\small
\textbf{Interpretation:} EU dictionary emphasises
\emph{import-side} terminology (\texttt{import\_duty}, \texttt{import\_restrictions});
US dictionary emphasises \emph{protectionist} and \emph{retaliatory} framing
consistent with the US trade policy debate.

\end{frame}

% ═══════════════════════════════════════════════════════
\section{Summary}
% ═══════════════════════════════════════════════════════

\begin{frame}{Summary \& Next Steps}

\begin{columns}[T]
\column{0.48\textwidth}
\begin{block}{Done}
\begin{itemize}\small
  \item GTM-8 baseline index (US, 1996--2025)
  \item Identified seed weaknesses $\to$ revised seeds for 4 of 8 topics
  \item Index 8 Revised: cleaner sub-topic separation,
        richer export-controls and financial-coercion signal
  \item Robustness check: article- vs.\ day-weighted ($r > 0.95$)
  \item External validation vs.\ GPR: $r \approx 0.13$;
        GEP Granger-causes GPR (lag 2+)
  \item Europe GEP index (independently trained GTM + dictionary)
  \item US--EU comparison with EPU-style normalisation
\end{itemize}
\end{block}

\column{0.48\textwidth}
\begin{block}{Next steps}
\begin{itemize}\small
  \item Predictive regressions: GEP on stock-market indices (US \& EU)
  \item Structural break tests around 2018 trade war, 2022 Ukraine
  \item Regime-switching model (calm vs.\ high-pressure episodes)
  \item Finalise thesis write-up and methodology chapter
\end{itemize}
\end{block}
\end{columns}

\bigskip
\begin{alertblock}{Key finding so far}
  The GEP index is \emph{distinct} from geopolitical risk (low correlation
  with GPRD), suggesting it captures a novel and complementary signal
  for financial markets.
\end{alertblock}

\end{frame}

% ─────────────────────────────────────────────
\end{document}
